\input{../../v3_rp/preamble}
\graphicspath{{../../}}
\makeatletter\def\input@path{{./}{../../}{../../v3_rp/}}\makeatother
\title{\bfseries Who Profits When a New Kind of Offering Appears?\\
\large Brief: the argument and the evidence}
\author{Eric Silver}
\date{October 2026}

\begin{document}
\maketitle

\section*{The question}

When a new kind of offering appears, the value it creates can go to the
firm that brought it, to the firms that follow, or to established firms
that already hold what is needed to sell it. The appropriability account
says which: the creator keeps the value where the offering is hard to
imitate and where the creator holds the complementary assets, such as
distribution, manufacturing and know-how; otherwise imitators and the
holders of those assets take it \citep{Teece1986}. The entry-order
literature argues the same question from named markets: pioneers keep
the markets they open \citep{RobinsonFornell1985, LiebermanMontgomery1988},
most pioneers fail and later entrants lead \citep{GolderTellis1993}, and
the answer depends on the pace of change \citep{SuarezLanzolla2007}.
These claims rest mostly on firms that survived to be studied.

\section*{Record, measure and outcome}

The US trademark record covers every entrant, including those that
failed. Each of 3.1 million registrations issued 2002--2018 carries a
dated description of what it sells. Between the fifth and sixth year
after registration the owner must swear that the mark is still in
commercial use on those goods, or the registration is cancelled; the
outcome is \emph{continued commercial use at five years}. Each
description is scored on two positions within its industry
\citep{SilverCorpus2026}: \emph{atypicality}, how unusual its language
is, which largely records language brought in from other industries, and
\emph{lead}, whether the industry's later filings moved toward it. Effects
below compare the top and bottom of each measure within an industry and
year, in points of survival.

\section*{What the record shows}

\emph{1. New offerings create value; arriving early loses it.} The most
unusual filing in its class and year is still in use 5.4 points more
often than the most typical, and renewed at year ten 4.3 points more
often. The most leading filing is in use 2.2 points less often than the
most lagging. No later outcome compensates: leading firms are no more
likely to raise private money, to list, or to be bought.

\emph{2. The creator keeps the value where it can be protected.} Owners
who patent keep 7.6 points of the advantage of unusual offerings; owners
who do not, 3.7. Imitation erodes it: the advantage falls from about 6
points to about 3 between industries whose filing volume grows slowly and
those where it grows fast, and from 6.2 to 2.8 as rivals in the same
theme concentrate in place. Patents do nothing for the cost of being
early.

\emph{3. The cost of being early belongs to the product, and part of it
is crowding.} Comparing each firm's leading marks with its own lagging
marks, two-thirds of the cost remains ($-0.82$ points against $-1.22$
without the comparison); it is not mainly a matter of weak firms chasing
trends. Entry by new firms into the same theme afterwards accounts for
about a fifth of it. Within new markets, pioneers and their early
followers fail alike.

\emph{4. The value goes to those who hold the complementary assets.} In
the surges of new themes, established firms' registrations survive most
often (54\% against 45\% for other entrants) and end up holding more than
their share of the marks still in use (19\% from 17\% of filings). The
cost of being early falls on entrants without those assets: self-filed
registrations lose survival steadily as lead rises, counsel-represented
ones only in the most leading tenth; platform businesses pay about 6
points; and the cost is larger in segments made up mostly of first-time
filers.

\emph{5. Industry structure predicts which products survive.} Measures of
the five forces \citep{Porter1980} built from the record explain half of
the variation in survival across 60 product segments. Blind ratings of
the same segments by three language models agree: rivalry and
substitutes go with lower survival (correlations about $-0.55$ for each
rater), and where a business depends on platforms and suppliers, being
early costs more.

\section*{What the record does not support}

Failed pioneers do not leave value that survivors collect. Later filings
that use language first introduced by filings that went on to fail
survive less often than their peers, not more.

\section*{The answer}

A new offering's value goes to its creator where the offering can be
protected and the creator holds the assets needed to sell it; elsewhere
it goes to followers and to established firms. The claim that pioneers
keep their markets and the claim that followers take them each describe
one side of that condition.

\section*{Limits}

Lead and atypicality measure the position of a filing's language, not
its market share. Continued use records whether a product was still
sold, not its profit. The estimates are associations within industry and
year; the comparisons within firms rule out stable differences between
firms, not choices that vary across a firm's own products.

\begin{thebibliography}{99}

\bibitem[Golder \& Tellis(1993)]{GolderTellis1993}
Golder, P.~N., \& Tellis, G.~J. (1993). Pioneer advantage: Marketing logic or
marketing legend? \emph{Journal of Marketing Research}, 30(2), 158--170.

\bibitem[Lieberman \& Montgomery(1988)]{LiebermanMontgomery1988}
Lieberman, M.~B., \& Montgomery, D.~B. (1988). First-mover advantages.
\emph{Strategic Management Journal}, 9(S1), 41--58.

\bibitem[Porter(1980)]{Porter1980}
Porter, M.~E. (1980). \emph{Competitive strategy: Techniques for analyzing industries and
competitors}. Free Press.

\bibitem[Robinson \& Fornell(1985)]{RobinsonFornell1985}
Robinson, W.~T., \& Fornell, C. (1985). Sources of market pioneer advantages
in consumer goods industries. \emph{Journal of Marketing Research}, 22(3),
305--317.

\bibitem[Silver(2026)]{SilverCorpus2026}
Silver, E. (2026). An event-dated corpus of US trademark prosecution and
a two-sided measure of vocabulary position. SSRN Working Paper 7520758.
\url{https://ssrn.com/abstract=7520758}.

\bibitem[Su\'arez \& Lanzolla(2007)]{SuarezLanzolla2007}
Su\'arez, F.~F., \& Lanzolla, G. (2007). The role of environmental
dynamics in building a first mover advantage theory. \emph{Academy of
Management Review}, 32(2), 377--392.

\bibitem[Teece(1986)]{Teece1986}
Teece, D.~J. (1986). Profiting from technological innovation:
Implications for integration, collaboration, licensing and public
policy. \emph{Research Policy}, 15(6), 285--305.

\end{thebibliography}

\end{document}
